% ===================================================================
\section{Tarea 6: riesgo crediticio multicapa}

\diapo{13}{De entidades a clientes de un banco}{1 min}
\begin{decir}
Hasta aquí las entidades eran abstractas. La tarea 6 lleva el modelo a la cartera de un
banco: los nodos son personas y empresas clientes del banco.

No existen datos públicos con la red de pagos de los clientes de un banco. Por eso
generamos una cartera sintética con un mecanismo conocido y verificamos que el método lo
encuentra. Es el mismo principio con que se valida un instrumento de medida: primero se
mide algo cuyo valor se conoce.

La cartera se construye para producir exactamente lo que un banco tendría: los atributos
de cada cliente y el registro de sus transacciones. El banco no ve el modelo interno; solo
ve los pagos.

Las relaciones están en tres capas. Nómina: la empresa paga a su empleado cada quince o
treinta días; es la dependencia más fuerte. Proveedores: una empresa paga a otra según
sus facturas. Transferencias: pagos entre personas, irregulares y de menor monto.
\end{decir}
\begin{terminos}
\begin{description}
\item[Cartera.] Conjunto de clientes con crédito en un banco.
\item[Datos sintéticos.] Datos generados por simulación con reglas conocidas.
\item[Red multicapa.] Red con varios tipos de relación entre los mismos nodos; cada tipo
es una capa con su propia matriz $W^\ell$ (Boccaletti et al., 2014).
\item[F y P.] Firma (empresa) y persona.
\item[Capa N, S, T.] Nómina, proveedores (\emph{suppliers}) y transferencias. En el
registro aparece además la capa C: cuotas de crédito pagadas al banco.
\end{description}
\end{terminos}
\begin{preguntas}
\pr{¿Qué valor tiene un resultado con datos sintéticos?}{Demuestra que, si el mecanismo
existe, el método lo detecta con la información que un banco tiene, y que no detecta
nada cuando el mecanismo no existe. No demuestra que el mecanismo exista en una cartera
real; eso requiere datos reales y es el siguiente paso.}
\pr{¿Hay trabajos similares?}{Letizia y Lillo (2019) usan redes de pagos entre empresas
para calificación crediticia; Tobback et al.\ (2017) y Óskarsdóttir y Bravo (2021) usan
redes relacionales y multicapa para predecir impago. Ninguno usa la sincronización de los
ciclos de pago.}
\end{preguntas}

\diapo{14}{Cartera sintética}{1.5 min}
\begin{decir}
La cartera de referencia tiene 2000 personas y 100 empresas. El 55~\% de las personas
son informales: no reciben nómina de una empresa del banco. La cifra es del orden de la
informalidad laboral reportada por el DANE.

El tamaño de las empresas sigue una distribución de Pareto: pocas empresas grandes y
muchas pequeñas. Cada empresa paga nómina mensual o quincenal. Los empleados formales
heredan el ritmo de su empleador con una variación del 8~\%; los informales tienen
ritmos más dispersos, con una variación del 25~\%.

La figura es un ejemplo reducido para ver las capas: los cuadrados son empresas y los
puntos son personas.

Cada cliente tiene una probabilidad de impago verdadera. Se calcula con una función
logística de sus variables clásicas más un componente oculto que el banco no ve. Ese
componente se revela en parte a través del historial de mora, como en la realidad. Las
columnas con la verdad del modelo empiezan por \texttt{truth\_} y el código impide que
cualquier modelo las use: si alguien lo intenta, el programa se detiene con un error.
\end{decir}
\begin{terminos}
\begin{description}
\item[Informalidad.] Trabajo sin vínculo laboral formal. En el modelo: persona sin
empleador dentro de la red del banco.
\item[DANE.] Departamento Administrativo Nacional de Estadística de Colombia.
\item[Distribución de Pareto.] Distribución de cola pesada, $P(X>x)\propto x^{-a}$. Aquí
$a=1.6$ para el tamaño de las empresas.
\item[Coeficiente de variación.] Desviación estándar dividida por la media. 8~\% y
25~\% son los coeficientes de variación de la frecuencia de pago.
\item[PD verdadera.] $\mathrm{logit}(\mathrm{PD}_i)=b_\text{tipo}+\sum_kc_kz_k(i)+
\sigma_uu_i$.
\item[$\mathrm{logit}(p)$.] $\ln\frac{p}{1-p}$, el logaritmo de la razón de
probabilidades.
\item[$z_k(i)$.] Variables clásicas estandarizadas (media 0, desviación 1) dentro de cada
tipo de cliente.
\item[$c_k$.] Pesos: endeudamiento $+0.7$, ingreso $-0.5$, antigüedad $-0.3$, edad
$-0.2$, formalidad $-0.5$.
\item[$u_i$, $\sigma_u$.] Componente latente $u_i\sim N(0,1)$ y su escala.
\item[Mora observada.] Número de moras en 12 meses $\sim\mathrm{Poisson}(m\,e^{au_i})$:
quien tiene $u_i$ alto tiende a tener más moras.
\item[$b_\text{tipo}$.] Intercepto calibrado por bisección para que la PD media sea 5~\%
en personas y 3~\% en empresas.
\item[Bisección.] Método numérico que encuentra la raíz de una función partiendo un
intervalo por la mitad repetidamente.
\item[SMMLV.] Salario mínimo mensual legal vigente. Unidad de los ingresos y montos.
\end{description}
\end{terminos}

\diapo{15}{Modelo dinámico multicapa}{1.5 min}
\begin{decir}
Esta es la versión del modelo para la cartera. Tiene la misma estructura que la ecuación
general con tres cambios.

Primero, el tiempo se mide en meses y el déficit como fracción del ingreso total del
cliente. Un déficit de 0.5 significa medio mes de ingreso. El ingreso total incluye lo que
llega de fuera de la red del banco.

Segundo, el contagio y el acople de fase tienen un valor por capa. La nómina transmite
más que las transferencias. Además, cada cliente ajusta su ritmo según la participación
de cada pagador en su ingreso: el empleador pesa mucho más que un amigo que le transfiere
dinero de vez en cuando.

Tercero, hay ruido aleatorio que representa choques individuales de liquidez, y hay
impagos idiosincráticos: cada cliente puede caer por razones propias, con una tasa que
corresponde a su probabilidad de impago.

El colchón depende del riesgo: un cliente con PD alta tiene un colchón más pequeño.

Con una sola capa y sin ruido, este modelo coincide con el de las tareas 1 a 4. Lo
verificamos numéricamente.
\end{decir}
\begin{terminos}
\begin{description}
\item[$\ell$.] Índice de capa: N, S o T.
\item[$W^\ell_{ij}$.] Flujo mensual de $j$ a $i$ en la capa $\ell$ [SMMLV/mes].
\item[$s_i$.] Ingreso total de $i$, incluido el externo a la red [SMMLV/mes].
\item[$k_i$.] Ingreso que llega por la red, $\sum_\ell\sum_jW^\ell_{ij}$.
\item[$\beta_\ell$.] Contagio por capa [1/mes]: N 0.5, S 0.4, T 0.3.
\item[$K_\ell$.] Acople de fase por capa [rad/mes]: N 0.8, S 1.0, T 0.3.
\item[$W^\ell_{ij}/k_i$.] Participación del pagador $j$ en el ingreso de red de $i$
(normalización por fila).
\item[$\gamma\frac{k_i}{s_i}(1-r_i)$.] Déficit por desfase, proporcional a la fracción del
ingreso que llega por la red. $\gamma=0.5$/mes.
\item[$\sigma\xi_i(t)$.] Ruido blanco de intensidad $\sigma=0.3$. Se integra con el
método de Euler–Maruyama.
\item[Euler–Maruyama.] Integrador de ecuaciones estocásticas:
$x_{n+1}=x_n+f(x_n)\Delta t+\sigma\sqrt{\Delta t}\,N(0,1)$.
\item[Colchón $x_i^{(c)}$.] $c\,(\mathrm{PD}_\text{ref}/\mathrm{PD}_i)^\kappa\varepsilon_i$
con $c=1.1$, $\mathrm{PD}_\text{ref}=0.04$, $\kappa=0.35$ y $\varepsilon_i$ lognormal.
\item[Impago idiosincrático.] Instante $\tau_i\sim\mathrm{Exp}(\lambda_i)$ con
$\lambda_i=-\ln(1-\mathrm{PD}_i)/12$ por mes, de modo que $P(\tau_i<12)=\mathrm{PD}_i$.
\item[Pérdida $\lambda=1$.] Perder un pagador que aporta todo el ingreso suma un mes de
ingreso al déficit.
\item[Matrices dispersas.] Almacenan solo las entradas distintas de cero; permiten
carteras grandes.
\end{description}
\end{terminos}
\begin{preguntas}
\pr{¿De dónde salen los valores de los parámetros?}{Son ilustrativos, no calibrados. Se
eligieron para que la tasa de impago y $R$ estén en rangos plausibles (impago cercano al
10~\%, $R$ entre 0.09 y 0.17). La diapositiva 22 muestra cómo cambia el resultado con
$\gamma$.}
\end{preguntas}

\diapo{16}{Del modelo a lo que observa el banco}{1 min}
\begin{decir}
Ahora el paso clave: pasar del modelo a lo que el banco realmente ve.

Cada vez que la fase de un pagador completa un ciclo, el pagador ejecuta sus pagos: a
sus empleados, a sus proveedores, y su cuota al banco. Agregamos ruido en la fecha y en el
monto. Quien cae en impago deja de pagar. El resultado es una tabla con fecha, pagador,
receptor, capa y monto: la misma estructura de un registro bancario.

Dividimos el tiempo en dos ventanas. De 0 a 6 meses, observamos: solo esas transacciones
alimentan las variables. De 6 a 18 meses, medimos si el cliente cae en impago; esa es la
etiqueta. Quien cae durante la observación se excluye.

Con las transacciones estimamos tres cosas: la red $\hat W$, la fase de cada cliente
$\hat\theta$ y su coherencia con los pagadores $\hat r$. La fase se estima por eventos:
entre dos pagos consecutivos, la fase avanza de manera lineal una vuelta completa.

La figura compara la fase real con la estimada para un empleado formal y un informal. La
estimación es buena: el error mediano en la red es menor al 5~\% y la correlación entre la
coherencia estimada y la real supera 0.8.
\end{decir}
\begin{terminos}
\begin{description}
\item[Registro de transacciones.] Tabla \texttt{(t, src, dst, layer, amount)}: instante,
pagador, receptor, capa y monto.
\item[Pago por ciclo.] El pagador $j$ paga a $i$ un monto $W_{ij}\cdot2\pi/\omega_j$ por
ciclo, de modo que el flujo mensual medio es $W_{ij}$.
\item[Ciclo observado.] Conjunto de pagos de una misma entidad separados por menos de
medio día; incluye la cuota al banco.
\item[Ruido de fecha.] $\pm1$ día, compartido por todos los pagos del mismo ciclo.
\item[Ventana de observación.] Meses $[0,6)$.
\item[Ventana de predicción.] Meses 6 a 18. Etiqueta $y=1$ si el cliente cae en ese
periodo.
\item[$\hat W_{ij}$.] Monto total pagado por $j$ a $i$ en la observación, dividido por
el periodo activo de $j$: $\min(t_\text{obs},\ \text{último ciclo}+\text{periodo medio})$.
Así, si $j$ dejó de pagar a mitad de la ventana, $\hat W_{ij}$ refleja lo que $i$ recibía
mientras $j$ pagaba, que es lo que perdió.
\item[Fase por eventos.] Pikovsky, Rosenblum y Kurths (2001):
$\hat\theta(t)=2\pi\bigl(k+\frac{t-t_k}{t_{k+1}-t_k}\bigr)$, donde $t_k$ es el instante del
$k$-ésimo ciclo de pagos.
\item[$\hat r_i(t)$.] $\frac1{\hat k_i}\sum_j\hat W_{ij}\cos(\hat\theta_j-\hat\theta_i)$,
calculada en una malla de 0.05 meses.
\item[Fuga de información.] Uso, al entrenar, de información que no estaría disponible al
predecir. El código rechaza transacciones posteriores a la observación.
\end{description}
\end{terminos}
\begin{preguntas}
\pr{¿Cómo se estima la fase de alguien que recibe pero casi no paga?}{Su fase se estima
con sus propios ciclos de pago, incluida la cuota al banco. Si no tiene ingreso por la
red, no hay desincronización medible y se le asigna $\hat r=1$.}
\end{preguntas}

\diapo{17}{Variables en tres niveles y protocolo}{1 min}
\begin{decir}
Con lo observado construimos variables en tres niveles, cada uno contiene al anterior.

El nivel 1 es el score clásico: ingreso, endeudamiento, edad, antigüedad, productos,
moras y formalidad.

El nivel 2 agrega la red estática: cuánto del ingreso llega por la red y por cuál capa,
qué tan concentrado está en pocos pagadores, la centralidad del cliente, y la salud de su
empleador y de sus pagadores. Esto corresponde a lo que ya explora la literatura.

El nivel 3 agrega la dinámica, que es el aporte del proyecto: la coherencia media y su
variabilidad, el tiempo que el cliente pasa desincronizado, el desfase con sus pagadores,
la diferencia de ritmo, la irregularidad de los ciclos y el ingreso de pagadores que
dejaron de pagar.

El protocolo tiene dos regímenes. Con red, los acoples están encendidos. En el régimen
nulo, apagamos el contagio, el desfase y la pérdida, pero mantenemos los mismos clientes,
la misma red y los mismos números aleatorios. Generamos 20 carteras por régimen y
validamos con GroupKFold por cartera. Comparamos los niveles por pares dentro de cada
partición y reportamos intervalos de confianza del 95~\%.
\end{decir}
\begin{terminos}
\begin{description}
\item[Concentración (HHI).] Índice de Herfindahl–Hirschman del ingreso,
$\sum_j(\hat W_{ij}/\hat k_i)^2$: vale 1 si todo viene de un pagador.
\item[PageRank.] Centralidad que asigna más importancia a quien recibe de entidades
importantes; factor de amortiguación 0.85. Se reescala para que la media sea 1.
\item[Salud del empleador y de los pagadores.] Tamaño, endeudamiento y moras del
empleador; moras y endeudamiento medios de los pagadores, ponderados por lo que pagan.
\item[$\hat r$ medio y desviación.] Promedio y desviación estándar de $\hat r_i(t)$ en la
observación.
\item[Tiempo desincronizado.] Fracción del tiempo con $\hat r_i<0.5$.
\item[Desfase.] Media en el tiempo de $|\arg\sum_j\hat W_{ij}e^{i(\hat\theta_j-
\hat\theta_i)}|$: el ángulo del desfase promedio con los pagadores, ponderado por lo que
paga cada uno, entre $0$ y $\pi$.
\item[Desajuste de ritmo.] $|\hat\omega_i-\bar\omega_\text{pagadores}|/2\pi$, en ciclos
por mes.
\item[Irregularidad de ciclos.] Coeficiente de variación de los intervalos entre ciclos,
propio y de los pagadores.
\item[Ingreso perdido.] Fracción del ingreso de red que venía de pagadores que dejaron de
pagar durante la observación.
\item[Régimen nulo.] $\beta_\ell=\gamma=\lambda=0$. Las fases y los pagos existen, pero no
crean déficit.
\item[Mismos sorteos.] La semilla de cada cartera genera cuatro flujos aleatorios
independientes (umbrales, impagos idiosincráticos, ruido y transacciones) con
\texttt{SeedSequence}. Los dos regímenes usan exactamente los mismos.
\item[Diferencia pareada.] Diferencia de métricas entre dos niveles evaluados en la misma
partición. Elimina la variación debida a qué carteras quedaron en cada partición.
\item[Intervalo de confianza 95~\%.] $\bar d\pm t_{0.975,4}\,s_d/\sqrt5$, con $\bar d$ la
media y $s_d$ la desviación de las 5 diferencias, y $t_{0.975,4}=2.78$ el cuantil de la
distribución $t$ de Student con 4 grados de libertad.
\end{description}
\end{terminos}
\begin{preguntas}
\pr{¿Por qué comparar contra la red estática y no solo contra el score?}{Porque superar
al score con información de red ya está documentado. La pregunta es si la dinámica añade
algo que la estructura no tiene.}
\end{preguntas}
